\documentclass[10pt,a4paper]{amsart}
\usepackage[margin=19mm]{geometry}
\usepackage{amsmath,amssymb,mathtools}
\usepackage[scr=rsfs,cal=euler]{mathalfa}
\usepackage{xfrac}
\usepackage[unicode,hidelinks,bookmarks=false]{hyperref}
\newcommand{\ie}{i.e.\ }
\newcommand{\cf}{cf.\ }
\newcommand{\resp}{respectively\ }
\newcommand{\Subsection}{\subsection}
\providecommand{\nobreakdash}{\nobreak\hbox{-}}
\setlength{\emergencystretch}{2em}
\multlinegap0pt
\usepackage{amssymb}
\usepackage{mathtools}
\usepackage{tikz}
\newtheorem{lemma}{Lemma}
\newcommand{\CC}{\mathbb C}
\newcommand{\ZZ}{\mathbb Z}
\title{\normalfont\LARGE A Complex Structure on $S^2\times S^4$}
\author{Shengtao Guo, Ethan X. Fang, Junwei Lu}
\date{Source: arXiv:2609.05504v1 (CC BY 4.0)}
\begin{document}
\maketitle

\noindent\emph{Source and licence.} \normalfont\LARGE A Complex Structure on $S^2\times S^4$. Version arXiv:2609.05504v1 (CC BY 4.0), distributed by arXiv under the Creative Commons Attribution 4.0 International licence, http://creativecommons.org/licenses/by/4.0/, as recorded in the arXiv licence metadata for this exact version and shown on the abstract page https://arxiv.org/abs/2609.05504v1. Full licence notice: \texttt{licenses/arxiv-2609.05504-CC-BY-4.0.txt}.\par\smallskip

\noindent\textit{Editorial scope.} Counted unit: the article's lemma lem:an-finite-index-cusp (source0.tex lines 861--883). The article's complete original proof of it is delivered with it (source0.tex lines 885--974, one \texttt{proof} environment of 90 lines). No mathematics is written, repaired or re-derived: the statement and the proof are the article's own lines, in the article's own notation, and no retained line is substituted or re-flowed.  No further source-local context is needed: every cross-reference of every retained line resolves inside the two delivered spans.  Nothing was omitted from any retained span, so the package carries no omission record.  No retained line contains a citation, so the wrapper carries no bibliography and no bibliography entry.  No retained line contains a figure environment, an image-inclusion command or drawing code, so no image is delivered and the package declares no asset dependency.  Editorial formatting only, in addition: an AMS \texttt{amsart} preamble that restates the article's own declarations and macros used by the retained lines, namely the environments \texttt{lemma} on separate counters, together with the standard packages those lines need.  The retained environments print in this document as Lemma 1 in order of appearance, whereas the article prints the same items with the numbering of its own full document.  The complete native source of the article as distributed in the arXiv e-print for this exact version is bundled unchanged as \texttt{sources/209441/source0.tex}.
\begin{lemma}[Toric compactification at the cusp]\label{lem:an-finite-index-cusp}
Let $\mathcal T$ be the periodic triangulation just defined, and let
$\Sigma$ be the fan in
$K_{\mathbb R}\oplus\mathbb R$ consisting of the cones over its faces
at height one.  Let $B_0:\ZZ^2\to K$ be injective with finite-index
image, and let $C(t)$ be a holomorphic $2$-by-$2$ matrix near zero.
The height character $t:=\chi^{(0,1)}$ extends to a holomorphic map
$t:X_\Sigma\to\CC$.  On the dense torus, write $x$ for the coordinate
in the two-dimensional torus with cocharacter lattice $K$, and define
\begin{equation}\label{eq:an-toric-action}
 \Psi_\lambda(x,t)
 =\bigl(e^{2\pi iC(t)\lambda}t^{B_0\lambda}x,t\bigr),
 \qquad \lambda\in\ZZ^2,
\end{equation}
where exponentials and powers are taken coordinatewise in the chosen basis
of $K$.  The maps extend to biholomorphisms.  For sufficiently
small $\epsilon>0$, their action on
$X_{\Sigma,\epsilon}:=\{|t|<\epsilon\}$ is free and properly
discontinuous, and
$X_{\Sigma,\epsilon}/\ZZ^2\to\{|t|<\epsilon\}$ is a proper map from a
smooth complex threefold.  Its central fibre is a reduced divisor with
normal crossings.
\end{lemma}

\begin{proof}
The integral shear $(y,r)\mapsto(y+rB_0\lambda,r)$
preserves $\Sigma$: at height one it translates $\mathcal T$ by the
lattice element $B_0\lambda$.  The induced toric map, composed with the
torus translation in \eqref{eq:an-toric-action}, gives the desired
extension.  The shear is the identity on the fibre cocharacter lattice and
fixes $t$, so it commutes with the torus factor.  These extensions form a
$\ZZ^2$-action because $\lambda\mapsto C(t)\lambda$ is additive.

For $t\ne0$, put $y=\log|x|/\log|t|$, coordinatewise; this is the
normalized logarithmic coordinate.  Then
\begin{equation}\label{eq:an-log-coordinate-translation}
 y\longmapsto y+B_t\lambda,
 \qquad
 B_t=B_0-\frac{2\pi\Im C(t)}{\log|t|}.
\end{equation}
Since $B_0$ is invertible over $\mathbb R$, shrinking $\epsilon$
gives a constant $a>0$ such that
\begin{equation}\label{eq:an-singular-value-bound}
 \|B_t\lambda\|\geq a\|\lambda\|
 \qquad(0<|t|<\epsilon,\ \lambda\in\ZZ^2).
\end{equation}
Hence the action is free on the dense torus.  A boundary
orbit is indexed by a bounded face $P$ of $\mathcal T$, and
$\Psi_\lambda$ carries it to the orbit indexed by
$P+B_0\lambda$.  No nonzero translation stabilizes a bounded face, so
the action is free on the boundary as well.

For proper discontinuity, let
$\sigma=\{v_0,v_1,v_2\}$.  Unimodularity identifies its chart with
$\CC^3$, with coordinates $z_0,z_1,z_2$ satisfying
\begin{equation}\label{eq:an-local-t-equation}
 t=z_0z_1z_2.
\end{equation}
If $\theta_i^\sigma$ are the barycentric coordinates of $\sigma$, then
on the dense torus
\begin{equation}\label{eq:an-barycentric-chart}
 \log|z_i|=\log|t|\,\theta_i^\sigma(y).
\end{equation}
Write
$U_\sigma(S)=\{|z_i|<S\text{ for }i=0,1,2\}
\cap X_{\Sigma,\epsilon}$.
Every point of $X_{\Sigma,\epsilon}$ lies in the closed unit polydisc of
some chart.  For a torus point, choose a triangle containing $y$.  For a
point on the orbit of a face $P$, choose a triangle containing $P$ whose
quotient cone contains the negative logarithmic coordinate normal to that
orbit.  Equation \eqref{eq:an-barycentric-chart} gives the assertion in
both cases.

If
$\Psi_\lambda U_\sigma(2)\cap U_{\sigma'}(2)\ne\varnothing$, density
gives such an intersection with $t\ne0$.  Formula
\eqref{eq:an-barycentric-chart} places the two normalized logarithmic
coordinates in
fixed bounded enlargements of $\sigma$ and $\sigma'$; their difference
is $B_t\lambda$.  The lower bound
\eqref{eq:an-singular-value-bound} therefore bounds $\lambda$ for each
fixed pair $(\sigma,\sigma')$.  The sets $U_\sigma(2)$ form an open
cover.  Given two compact sets, choose
finite subcovers of each by these sets and apply the bound to the finitely
many resulting pairs of triangles.  Only finitely many group elements can
move one compact set to meet the other, which proves proper discontinuity.

For properness over the disc, first suppose $t\ne0$ and choose
$\lambda$ so that
$B_t^{-1}y+\lambda\in[-1/2,1/2]^2$.
Uniform upper bounds on $B_t$ then put the translated logarithmic coordinate
in a fixed ball.  At $t=0$, use the finite set
$K/B_0\ZZ^2$: translating a face makes one of its vertices belong to a
fixed set of coset representatives.  Local finiteness of $\mathcal T$
then leaves only finitely many adjacent triangles.  Consequently, for each
$0<\epsilon_1<\epsilon$, every orbit over
$|t|\leq\epsilon_1$ meets
\begin{equation*}
 \mathcal K_0=
 \bigcup_{\sigma\in\mathcal F}
 \{|z_0|,|z_1|,|z_2|\leq1,\ |t|\leq\epsilon_1\},
\end{equation*}
where $\mathcal F$ is a finite family of triangles.  This is compact in
$X_{\Sigma,\epsilon}$, and its image is the full inverse image of the
closed subdisc.  The quotient map to the disc is therefore proper.

Each of the two triangle types has determinant $1$ or $-1$, so
$\Sigma$ is unimodular and $X_\Sigma$ is smooth.  The free quotient is
smooth.  The fan is countable, so $X_\Sigma$ has a countable toric atlas;
the quotient map is open, and the quotient is therefore second countable.
Since every ray generator has height one,
\eqref{eq:an-local-t-equation} gives a reduced normal crossings central
fibre.
\end{proof}

\end{document}
